\begin{lemma}[The sentence following (3.9)]
Let $\mathbb G$ be a sub-elliptic setting (\noderef{CarnotSetting}) with group law $*$ and
homogeneous norm $\|\cdot\|$, let $\Om\subset\mathbb R^{n}$ be a bounded open set, and let
$u,v:\mathbb R^{n}\to\mathbb R$ with $u$ locally semi-convex and $v$ locally semi-concave in
$\Om$, i.e. $-v$ locally semi-convex (\noderef{LocallySemiConvexOn}). Let $\delta>0$ and let
$h\in\mathbb R^{n}$ with $\|h\|<\delta$. Write $\Om_\delta$ for the inner domain
(\noderef{TranslatedDomain}) and $u_h:=u(\cdot*h)$. Then
\begin{enumerate}
\item $x*h\in\Om$ for every $x\in\Om_\delta$;
\item $u_h$ is locally semi-convex on $\Om_\delta$;
\item $-v$ is locally semi-convex on $\Om_\delta$, i.e. $v$ is locally semi-concave there;
\item for every $x\in\Om_\delta$ at which $u_h-v$ is maximal over $\Om_\delta$, there is a
  continuous linear functional $D$ which is the derivative of both $u_h$ and $v$ at $x$, and
  $u$ is differentiable at $x*h$.
\end{enumerate}
\end{lemma}
\begin{proof}
This is the sentence the paper places immediately after (3.9): "By (2.16), $x*h\in\Om$
whenever $x\in\Om_\delta$ and $\|h\|<\delta$. \dots The function $u_h$ is semi-convex on
compact subsets of $\Om_\delta$ (Section Section~2.3(4)), so at every point $x$ of
$\A_\delta(h)$ the functions $u_h$ and $v$ are differentiable, and $u$ is differentiable at
$x*h$."

\medskip\noindent\textbf{Step 0: a restriction principle for local semi-convexity.} We first
record the following, used twice below. \emph{If $w$ is locally semi-convex on a set $T$ and
$S\subseteq T$ is open, then $w$ is locally semi-convex on $S$.} Indeed, let $x\in S$. Then
$x\in T$, so by \noderef{LocallySemiConvexOn} there are $\rho>0$ and $\Lambda\ge0$ with
$B(x,\rho)\subseteq T$ and $y\mapsto w(y)+\frac\Lambda2|y|^{2}$ convex on the Euclidean ball
$B(x,\rho)$. Since $S$ is open there is $\rho'>0$ with $B(x,\rho')\subseteq S$. Put
$\rho'':=\min\{\rho,\rho'\}>0$. Then $B(x,\rho'')\subseteq B(x,\rho')\subseteq S$, and
$B(x,\rho'')\subseteq B(x,\rho)$ is a convex subset of $B(x,\rho)$, so the restriction of the
convex function $y\mapsto w(y)+\frac\Lambda2|y|^{2}$ to $B(x,\rho'')$ is convex. Thus
$\rho''$ and $\Lambda$ witness the defining condition of \noderef{LocallySemiConvexOn} for
$w$ on $S$ at $x$. As $x\in S$ was arbitrary, $w$ is locally semi-convex on $S$.

\medskip\noindent\textbf{Step 1: part (1).} Let $x\in\Om_\delta$. By associativity of $*$ and
$0*h=h$ (\noderef{CarnotSetting}),
\[
 d(x*h,x)=\big\|(-x)*(x*h)\big\|=\big\|\big((-x)*x\big)*h\big\|=\|0*h\|=\|h\| ,
\]
which is the exact identity (2.16) (kept inline by pm-0012). Since $\|h\|<\delta$,
the point $x*h$ lies in the $d$-ball of radius $\delta$ about $x$, and by item (6) of
\noderef{InnerDomainProperties} that ball is contained in $\Om$. Hence $x*h\in\Om$.

\medskip\noindent\textbf{Step 2: part (2).} Let $\Theta(y):=y*h$. By
\noderef{CarnotSetting} the group law is $C^{\infty}$, so $\Theta$, a partial application of
it, is $C^{\infty}$ and in particular $C^{2}$. The second clause of
\noderef{SemiConvexCompose}, applied with the set $\Om$, the function $u$ (locally
semi-convex on $\Om$ by hypothesis) and this $\Theta$, gives
\[
 u\circ\Theta=u_h\ \text{is locally semi-convex on}\ \Theta^{-1}(\Om)=\{y:\ y*h\in\Om\} .
\]
By Step 1, $\Om_\delta\subseteq\Theta^{-1}(\Om)$, and $\Om_\delta$ is open by item (3) of
\noderef{InnerDomainProperties}. Step 0 with $T:=\Theta^{-1}(\Om)$ and $S:=\Om_\delta$ now
gives that $u_h$ is locally semi-convex on $\Om_\delta$. This is the paper's appeal to
Section~2.3(4); note that we obtain local semi-convexity on the open set $\Om_\delta$
itself, which is the form \noderef{SemiConvexDiffAtMax} consumes, rather than the paper's
"on compact subsets".

\medskip\noindent\textbf{Step 3: part (3).} $-v$ is locally semi-convex on $\Om$ by
hypothesis, $\Om_\delta\subseteq\Om$ by item (1) of \noderef{InnerDomainProperties}, and
$\Om_\delta$ is open by item (3) of the same node. Step 0 with $T:=\Om$ and
$S:=\Om_\delta$ gives the claim.

\medskip\noindent\textbf{Step 4: part (4), the common derivative at a maximal point.} Let
$x\in\Om_\delta$ be a point at which $u_h-v$ is maximal over $\Om_\delta$. Apply the second
clause of \noderef{SemiConvexDiffAtMax} with the set $\Om_\delta$, the locally semi-convex
function $u_h$ (Step 2) and the locally semi-concave function $v$ (Step 3): since
$x\in\Om_\delta$ and $u_h-v$ attains its maximum over $\Om_\delta$ at $x$, there is a
continuous linear functional $D$ with
\[
 D\ \text{is the derivative of}\ u_h\ \text{at}\ x\quad\text{and}\quad
 D\ \text{is the derivative of}\ v\ \text{at}\ x . \tag{1}
\]

\medskip\noindent\textbf{Step 5: part (4), differentiability of $u$ at $x*h$.} Let
$\Theta'(z):=z*(-h)$, which is $C^{\infty}$ by \noderef{CarnotSetting} for the same reason as
$\Theta$. For every $z$, using associativity, $(-h)*h=0$ and $z*0=z$,
\[
 u_h\big(\Theta'(z)\big)=u\Big(\big(z*(-h)\big)*h\Big)=u\Big(z*\big((-h)*h\big)\Big)=u(z*0)=u(z),
\]
so $u=u_h\circ\Theta'$ identically on $\mathbb R^{n}$. Moreover
$\Theta'(x*h)=(x*h)*(-h)=x*\big(h*(-h)\big)=x*0=x$. Now $\Theta'$ is differentiable at
$x*h$, with derivative the continuous linear map $L:=D\Theta'(x*h)$, and $u_h$ is
differentiable at $\Theta'(x*h)=x$ with derivative $D$ by (1). By the chain rule,
$u=u_h\circ\Theta'$ is differentiable at $x*h$ with derivative $D\circ L$, a continuous
linear functional. This proves the second half of part (4), and completes the proof.
\end{proof}
